%% ma: L12-B08-0179
%% lop: 12
%% chuong: 2
%% bai: 8
%% dang: 12.8.53
%% loai: TN4
%% muc: VD
%% dap_an: "A"
%% nguon: "CĐ2.4.pdf giáo viên gửi 27/10/2026 (M07, Lê Minh Khôi – Nguyễn Anh Khôi, Chương II Chủ đề 4 Các bài toán về tâm tỉ cự), Câu 29"
%% kien_thuc: "Bài 8 (khoảng cách hai điểm, biểu thức tọa độ), Bài 6 (bình phương vô hướng, tâm tỉ cự), Bài 7 (hình chiếu lên trục, mặt phẳng tọa độ)"
%% trang_thai: cho-duyet
%% ly_do: "Dạng mới 12.8.53 đề xuất, chờ giáo viên duyệt; Đề gốc không có phương án đúng; đã sửa phương án A từ 116 thành 190"
%% ghi_chu: "Đề gốc không có phương án đúng (các phương án 116; 248; -184; 102, giá trị đúng là 190); đã sửa phương án A thành 190"
\begin{ex}
Trong không gian $Oxyz$, cho ba điểm $A(2;3;-1)$, $B(-2;1;-3)$, $C(-4;-5;1)$. Xét hai điểm $H$, $K$ thay đổi với $H$ thuộc mặt phẳng $(Oyz)$ và $K$ thuộc trục hoành sao cho biểu thức $T=3HA^2-HB^2-HC^2$ đạt giá trị nhỏ nhất và $Q=-3KA^2+2KB^2-KC^2$ đạt giá trị lớn nhất. Tính giá trị của biểu thức $V=2\cdot HK^2+\dfrac T2+2\cdot Q$.
\choice{\True $V=190$}{$V=248$}{$V=-184$}{$V=102$}
\loigiai{$I=3A-B-C=(12;13;-1)$; $T=HI^2+3IA^2-IB^2-IC^2=HI^2-328$, nhỏ nhất khi $H(0;13;-1)$, $T=144-328=-184$.

$J=\dfrac{-3A+2B-C}{-2}=(3;1;2)$; $Q=-2KJ^2-28$, lớn nhất khi $K(3;0;0)$, $Q=-38$.

$HK^2=179$, $V=358-92-76=190$.}
\end{ex}
